\cotefile{156-9}
\scanpage{18}
\mspage{14}
% archive sheet number: 17

des polygones et des figures, leurs $\delta\Phi$, $\partial\Phi$, ne sont pas changés, à un détail près, concernant \deleted{les} \inserted{définition et terminologie concernant les} ``\uline{sommets} \uline{intérieurs} \deleted{de $\Phi$} ou ``internes'' de $\Phi$, propres ou impropres. En plus des sommets internes pris précédemment (\deleted{des} $s\in\delta\Phi$ t.q. $\operatorname{ord}(s,\Phi)=2$), il faut rajouter les \deleted{$s\in\delta\Phi$} $s\in\partial I$ tels que $\operatorname{ord}(s,\unsure{\Phi|\mathbb{\Phi}})=1$, \inserted{i.e. $s\in\partial\Phi$.} (Car si $s\in\unsure{\delta I|\partial I}$, on a tjrs $\operatorname{ord}(s,\Phi)\in\{0,1\}$, à l'exclusion de \deleted{\unsure{2}} $2$, \deleted{et cela se} \unsure{le disant} pour $\forall\,\Phi$ ou pour $\forall\,\Phi$ de la forme $\{\deleted{I}\unsure{X|\alpha|\,}\}$, où $\alpha\in\unsure{M_1|\mathcal M_1}$, \unsure{cela caractérise} les pts de $\delta I$). On dit qu'un tel sommet est \deleted{\unsure{lac|non}} \uline{redondant} s'il est propre, ou encore si (comme pour les sommets redondants proprement dits) il n'est pas dans $\delta\Phi^*$; dans le cas inverse, il est dit sommet lacunaire. Ceci dit, on a
\[
\Phi^*_{\mathrm{is}}=\Phi_{\mathrm{lac}}\qquad \Phi^*_{\mathrm{lac}}=\Phi_{\mathrm{is}}
\]
\[
\Phi_{\mathrm{red}}=\delta\Phi-\delta\Phi^*
\]
\[
\delta\Phi=\underbrace{\Phi_{\mathrm{red}}\cup\Phi_{\mathrm{lac}}}_{\Phi_{\mathrm{int}}}\cup\Phi_{\mathrm{is}}\cup
\underbrace{\partial\Phi}_{
 \underbrace{\partial_0\Phi}_{\substack{\Phi_0\cap\delta\Phi\\\text{sommets-bord}\\\text{propres}}}
 \amalg
 \underbrace{\partial'_0\Phi}_{\text{somm.b. impropres}}}
\]
\margin{Écrire réunion \uline{disjointe}
$\delta\Phi=\delta\Phi=\Phi_{\deleted{\unsure{\mathrm{red}}}\inserted{\mathrm{red}}}\cup\Phi_{\deleted{\unsure{\mathrm{red}}}\mathrm{is}}\cup\Phi_{\mathrm{lac}}\cup(\unsure{\partial^*|\delta^*}\Phi\deleted{\unsure{\delta\Phi\cap}}$
où $\partial^*\Phi=\partial\Phi\setminus\delta\Phi\cap\partial I$ (sommets-bords \uline{non} \uline{extérieurs})}
réunions disjointes, sauf pour \unsure{le dernier terme|les derniers termes} $\partial\Phi$ et les deux premiers, formant $\Phi_{\mathrm{int}}$, car
\[
\partial\Phi\cap\Phi_{\mathrm{int}}=\partial I\cap\deleted{\unsure{\delta I\cap}}(\deleted{\unsure{\delta\Phi}}\Phi_0\setminus\Phi_{\mathrm{is}})
\]
les points de $\delta\Phi$ qui sont intérieurs sont les \deleted{points} qui sont dans $\partial I$, sans être des points isolés de $\Phi$.
